% !TEX program = pdflatex
\documentclass[aspectratio=169,10pt]{beamer}

\usepackage[T1]{fontenc}
\usepackage[utf8]{inputenc}
\usepackage[spanish]{babel}
\usepackage{booktabs}
\usepackage{graphicx}
\usepackage{hyperref}
\usepackage{xcolor}
\usepackage{tikz}
\usepackage{enumitem}
\setlist{leftmargin=*,itemsep=0.22em,topsep=0.15em}

\usetheme{Madrid}
\usecolortheme{default}
\setbeamertemplate{navigation symbols}{}
\setbeamertemplate{footline}{%
  \hfill{\scriptsize\color{black!60}\insertsectionhead\ \textbar\ \insertframenumber/\inserttotalframenumber}\hspace{1em}\vspace{0.35em}%
}

\definecolor{ieblue}{RGB}{0,51,102}
\definecolor{good}{RGB}{34,120,70}
\definecolor{bad}{RGB}{160,40,40}
\definecolor{soft}{RGB}{245,247,250}
\definecolor{accent}{RGB}{0,120,160}
\setbeamercolor{structure}{fg=ieblue}
\setbeamercolor{frametitle}{fg=ieblue,bg=soft}
\setbeamercolor{block title}{bg=ieblue,fg=white}
\setbeamercolor{block body}{bg=soft}
\setbeamerfont{frametitle}{size=\large,series=\bfseries}

% Shared TikZ helpers --- layouts seguros (sin solapes)
\usetikzlibrary{arrows.meta,positioning,shapes.geometric,calc}

\tikzset{
  box/.style={
    draw=ieblue, rounded corners=3pt, align=center,
    fill=soft, font=\scriptsize, inner sep=5pt,
    minimum height=8.5mm, text depth=0.3em
  },
  boxdark/.style={
    draw=ieblue, rounded corners=3pt, align=center,
    fill=ieblue, text=white, font=\scriptsize\bfseries, inner sep=5pt,
    minimum height=8.5mm, text depth=0.3em
  },
  arr/.style={-{Latex}, thick, ieblue},
  pill/.style={
    draw=ieblue, rounded corners=8pt, fill=white,
    font=\scriptsize\bfseries, inner sep=4pt
  }
}

% Pipeline horizontal con anchos fijos y hueco para flechas
\newcommand{\PipelineFive}[5]{%
\begin{center}
\begin{tikzpicture}
  \node[box, text width=1.85cm] (a) {#1};
  \node[box, text width=1.85cm, right=0.55cm of a] (b) {#2};
  \node[box, text width=1.85cm, right=0.55cm of b] (c) {#3};
  \node[box, text width=1.85cm, right=0.55cm of c] (d) {#4};
  \node[boxdark, text width=1.85cm, right=0.55cm of d] (e) {#5};
  \draw[arr] (a) -- (b);
  \draw[arr] (b) -- (c);
  \draw[arr] (c) -- (d);
  \draw[arr] (d) -- (e);
\end{tikzpicture}
\end{center}%
}

% Encaje Pedro / Mario con minipages (evita solapes TikZ)
\newcommand{\FitDual}[4]{%
\noindent
\begin{minipage}[t]{0.48\textwidth}
\centering
\begin{tikzpicture}\node[circle, draw=ieblue, fill=blue!12, minimum size=1cm, font=\small\bfseries]{Pedro};\end{tikzpicture}\\[0.25em]
\begin{beamercolorbox}[wd=\textwidth,rounded=true,sep=4pt]{block body}
\scriptsize #1\\[0.35em]\textbf{Afinidad: #3/10}
\end{beamercolorbox}
\end{minipage}\hfill
\begin{minipage}[t]{0.48\textwidth}
\centering
\begin{tikzpicture}\node[circle, draw=orange!70!black, fill=orange!18, minimum size=1cm, font=\small\bfseries]{Mario};\end{tikzpicture}\\[0.25em]
\begin{beamercolorbox}[wd=\textwidth,rounded=true,sep=4pt]{block body}
\scriptsize #2\\[0.35em]\textbf{Afinidad: #4/10}
\end{beamercolorbox}
\end{minipage}%
}

% Pros / contras con minipages
\newcommand{\ProCon}[2]{%
\noindent
\begin{minipage}[t]{0.48\textwidth}
\begin{beamercolorbox}[wd=\textwidth,rounded=true,sep=5pt]{block body}
\textcolor{good}{\textbf{Lo bueno}}\\[0.3em]
\scriptsize #1
\end{beamercolorbox}
\end{minipage}\hfill
\begin{minipage}[t]{0.48\textwidth}
\begin{beamercolorbox}[wd=\textwidth,rounded=true,sep=5pt]{block body}
\textcolor{bad}{\textbf{Lo malo}}\\[0.3em]
\scriptsize #2
\end{beamercolorbox}
\end{minipage}%
}

% Temario vs extensión
\newcommand{\MethodsSplit}[2]{%
\noindent
\begin{minipage}[t]{0.48\textwidth}
\begin{beamercolorbox}[wd=\textwidth,rounded=true,sep=5pt]{block body}
\textbf{Del curso}\\[0.25em]\scriptsize #1
\end{beamercolorbox}
\end{minipage}\hfill
\begin{minipage}[t]{0.48\textwidth}
\begin{beamercolorbox}[wd=\textwidth,rounded=true,sep=5pt]{block title}
\textbf{Más allá del temario}\\[0.15em]
\end{beamercolorbox}
\begin{beamercolorbox}[wd=\textwidth,rounded=true,sep=5pt]{block body}
\scriptsize #2
\end{beamercolorbox}
\end{minipage}%
}

\newcommand{\Verdict}[2]{%
\begin{center}
\begin{tikzpicture}
  \node[boxdark, text width=0.90\textwidth, align=center, font=\normalsize] {
    \textbf{\large #1}\\[0.35em]
    \normalfont\normalsize #2
  };
\end{tikzpicture}
\end{center}%
}

\newcommand{\ScoreBar}[3]{%
\noindent\begin{tikzpicture}
  \node[anchor=west, font=\scriptsize, text width=2.3cm] at (0,0) {#1};
  \fill[black!10] (2.4, -0.12) rectangle (6.6, 0.22);
  \fill[#3] (2.4, -0.12) rectangle ({2.4+#2*0.42}, 0.22);
  \node[font=\scriptsize\bfseries, anchor=west] at (6.75,0.05) {#2};
\end{tikzpicture}\par\vspace{0.1em}%
}

% Cabecera de opción sin solapes
\newcommand{\OptionHeader}[3]{%
\begin{center}
\begin{tikzpicture}
  \node[boxdark, text width=4.0cm] (a) {#1};
  \node[pill, right=0.5cm of a] (b) {#2};
\end{tikzpicture}
\end{center}
\vspace{0.35em}
\begin{block}{En una frase}
#3
\end{block}%
}


\title{Top 10 opciones de proyecto ML}
\subtitle{Pedro $\times$ Mario --- Decision deck visual y accesible}
\author{Research brief deep $\cdot$ Octubre 2026}
\date{Octubre 2026}
\institute{BAM $\cdot$ Machine Learning \& Communication Skills $\cdot$ IE University}

\begin{document}

\begin{frame}{Portada}
  \titlepage
\end{frame}

\begin{frame}{Para qué sirve este deck}
  \begin{columns}[T]
    \begin{column}{0.55\textwidth}
      \begin{itemize}
        \item Traduce el research técnico a una conversación de decisión.
        \item 10 opciones: 5 recomendadas + 5 condicionales.
        \item Cada opción se explica primero \textbf{para alguien de fuera del sector}.
        \item Objetivo final: elegir \textbf{1 principal + 1 reserva}.
      \end{itemize}
    \end{column}
    \begin{column}{0.42\textwidth}
      \begin{tikzpicture}
        \node[boxdark, text width=4.4cm] (a) {1. Entender el problema};
        \node[box, text width=4.4cm, below=0.28cm of a] (b) {2. Ver pros y contras};
        \node[box, text width=4.4cm, below=0.28cm of b] (c) {3. Marcar verde / ámbar / rojo};
        \node[boxdark, text width=4.4cm, below=0.28cm of c] (d) {4. Elegir principal + reserva};
        \draw[arr] (a)--(b); \draw[arr] (b)--(c); \draw[arr] (c)--(d);
      \end{tikzpicture}
    \end{column}
  \end{columns}
\end{frame}

\begin{frame}{Qué pedimos a cada idea (filtro del curso)}
  \begin{center}
  \begin{tikzpicture}
    \node[boxdark, text width=2.5cm] (d) {Datos reales\\manejables};
    \node[boxdark, text width=2.5cm, right=0.45cm of d] (m) {Métodos del\\temario (+extra)};
    \node[boxdark, text width=2.5cm, right=0.45cm of m] (a) {Ambición\\visible (25\%)};
    \node[boxdark, text width=2.5cm, right=0.45cm of a] (o) {Defendible\\en oral};
  \end{tikzpicture}
  \end{center}
  \vspace{0.5em}
  \begin{block}{Marco práctico}
    Python + scikit-learn + PyTorch \quad$\cdot$\quad $\sim$6 semanas \quad$\cdot$\quad informe 10--12 páginas + GitHub con historial real
  \end{block}
\end{frame}

\begin{frame}{Qué pedimos a cada idea (filtro nuestro)}
  \FitDual{Deeptech, física, robótica, espacio, quantum. Rechaza lo genérico. Quiere un proyecto que impresione. No busca el camino fácil.}{Finanzas y sistemas de recomendación. Quiere originalidad, buenos datos y comparación limpia. Dispuesto a leer papers y reproducir ideas.}{10}{8}
  \vspace{0.45em}
  \begin{block}{Vetos}
    \textbf{Mario:} nada de salud / biomedicina / omics. \quad
    \textbf{Pedro:} nada de proyecto genérico de estudiante.
  \end{block}
\end{frame}

\begin{frame}{Cómo leer cada bloque de opción}
  \begin{center}
  \begin{tikzpicture}
    \foreach \i/\t [count=\n] in {
      1/Portada,
      2/Sin jerga,
      3/Analogía,
      4/Qué hace el ML,
      5/Glosario,
      6/Pregunta,
      7/Pasos,
      8/Datos,
      9/Métodos,
      10/Encajes + veredicto} {
      \pgfmathsetmacro{\x}{mod(\n-1,5)}
      \pgfmathsetmacro{\y}{-floor((\n-1)/5)}
      \node[box, text width=2.15cm, minimum height=1.1cm] at ({2.35*\x},{1.35*\y}) {\scriptsize\textbf{\i}\\\t};
    }
  \end{tikzpicture}
  \end{center}
  \vspace{0.45em}
  Si no entendéis el problema en las diapositivas ``sin jerga'' / ``analogía'', esa opción aún no está lista para decidir.
\end{frame}

\begin{frame}{El top 10 de un vistazo}
  \begin{columns}[T]
    \begin{column}{0.48\textwidth}
      \begin{block}{Recomendadas}
        \scriptsize
        1. Evitar choques entre satélites (ESA)\\
        2. ¿Qué startups llegan a Series A?\\
        3. ¿Los modelos entienden o memorizan?\\
        4. Enseñar física a la red neuronal\\
        5. Ver incendios aunque haya nubes
      \end{block}
    \end{column}
    \begin{column}{0.48\textwidth}
      \begin{block}{Condicionales}
        \scriptsize
        6. Propiedades de materiales (Matbench)\\
        7. Perfil del founder $\rightarrow$ gran éxito\\
        8. ¿La IA entiende analogías de verdad?\\
        9. Cambios urbanos en dos fotos satélite\\
        10. Reutilizar experiencias de robots
      \end{block}
    \end{column}
  \end{columns}
\end{frame}

\section{Recomendadas}
\subsection{Opción 1 --- ESA Collision Avoidance}

\begin{frame}{1 · Portada --- Evitar choques entre satélites}
  \OptionHeader{Recomendada nº 1}{Espacio / ESA}{%
    Usamos avisos oficiales de la Agencia Espacial Europea para anticipar si dos objetos en órbita van a acercarse de forma peligrosa.
  }
\end{frame}

\begin{frame}{1 · Imagina esto (sin jerga)}
  \begin{columns}[T]
    \begin{column}{0.55\textwidth}
      \begin{itemize}
        \item Piensa en el espacio cercano a la Tierra como una \textbf{autopista muy llena}.
        \item Satélites, restos de cohetes y basura espacial circulan a miles de km/h.
        \item Cuando dos objetos se van a cruzar cerca, llega un \textbf{aviso} (como un parte meteorológico del riesgo).
        \item Ese aviso se actualiza varias veces: el riesgo sube o baja según llegan mejores medidas.
        \item Alguien tiene que decidir: \textbf{¿desvío el satélite o no?}
      \end{itemize}
    \end{column}
    \begin{column}{0.42\textwidth}
      \begin{tikzpicture}
        \draw[thick, ieblue] (0,0) circle (1.1);
        \node[font=\tiny] at (0,0) {Tierra};
        \draw[thick, accent] (0,0) circle (1.65);
        \fill[ieblue] (1.5,0.65) circle (0.1);
        \fill[bad] (1.85,0.4) circle (0.09);
        \node[font=\tiny, align=center] at (0.9,-1.95) {dos objetos\\se acercan};
      \end{tikzpicture}
    \end{column}
  \end{columns}
\end{frame}

\begin{frame}{1 · ¿Qué haría aquí el machine learning?}
  \begin{itemize}
    \item Hoy hay física + operadores humanos mirando muchos avisos.
    \item Nosotros preguntamos: \textbf{con los avisos tempranos}, ¿podemos estimar cómo de grave será el riesgo al final?
    \item Es parecido a: ``con el parte de las 10:00, ¿prevemos si la tormenta de las 18:00 será fuerte?''
    \item Si el modelo avisa bien y pronto, se pueden evitar maniobras innecesarias (gastan combustible) y no ignorar peligros reales.
  \end{itemize}
\end{frame}

\begin{frame}{1 · Palabras que vas a oír (glosario rápido)}
  \begin{description}[style=nextline,leftmargin=1.2cm,font=\bfseries\scriptsize]
    \item[CDM] Mensaje de ``casi choque'': una fila de datos con distancia, incertidumbre y riesgo estimado.
    \item[Evento] Toda la secuencia de avisos del mismo acercamiento (no una sola fila).
    \item[TCA] Momento del máximo acercamiento.
    \item[Desbalance] Casi todos los casos son ``poco peligro''; los graves son raros --- por eso no vale mirar solo ``aciertos totales''.
  \end{description}
\end{frame}

\begin{frame}{1 · Pregunta de investigación}
  \begin{block}{Pregunta principal}
    Si solo vemos parte de la secuencia de avisos (por ejemplo, hasta 2 días antes), ¿podemos predecir el riesgo del último aviso mejor que reglas simples?
  \end{block}
  \begin{block}{Pregunta secundaria}
    ¿Ayuda meter variables físicas (distancia, incertidumbre de posición) para reducir falsas alarmas?
  \end{block}
\end{frame}

\begin{frame}{1 · Cómo lo haríamos, paso a paso}
  \PipelineFive{Explorar\\datos}{Crear\\variables}{Modelos\\clásicos}{Redes +\\calibración}{Comparar\\con/sin física}
  \vspace{0.4em}
  \begin{itemize}
    \item Unidad de análisis: cada \textbf{acercamiento completo}.
    \item Éxito = detectar bien los casos raros y peligrosos, no acertar el 99\% fácil.
    \item Entregable oral: cuándo el modelo ayuda al operador y cuándo se equivoca.
  \end{itemize}
\end{frame}

\begin{frame}{1 · De dónde salen los datos}
  \begin{itemize}
    \item Dataset público de la \textbf{ESA} (reto Kelvins / Collision Avoidance).
    \item Miles de acercamientos reales (2015--2019) con muchas variables por aviso.
    \item Ventaja: no hace falta bajar imágenes de satélite; es tabla + series.
    \item Reto: hay papers de competición --- hay que aportar un \textbf{ángulo propio} (avisar antes, calibrar probabilidades, física).
  \end{itemize}
\end{frame}

\begin{frame}{1 · Métodos del curso y más allá}
  \MethodsSplit{Regresión / logística · Árboles / boosting · Validación cuidadosa · Ingeniería de variables}{Red secuencial pequeña · Calibración de probabilidades · Variables orbitales · Opcional: procesos gaussianos}
\end{frame}

\begin{frame}{1 · Por qué encaja con el curso}
  \begin{columns}[T]
    \begin{column}{0.55\textwidth}
      \begin{itemize}
        \item Datos reales y poco típicos de estudiante.
        \item Comparación de modelos natural.
        \item Limitaciones claras (rareza, incertidumbre).
        \item Oral defendible: física + métricas.
      \end{itemize}
    \end{column}
    \begin{column}{0.42\textwidth}
      \ScoreBar{Ambición}{9}{ieblue}
      \ScoreBar{Datos}{9}{accent}
      \ScoreBar{Factibilidad}{8}{good}
      \ScoreBar{Oral}{9}{ieblue}
    \end{column}
  \end{columns}
\end{frame}

\begin{frame}{1 · Por qué encaja con nosotros}
  \FitDual{Le encanta el espacio. Evita proyectos genéricos. Puede liderar la parte física y la historia SSA.}{Puede liderar métricas, casos raros, calibración y la comparación limpia de modelos.}{10}{7}
\end{frame}

\begin{frame}{1 · Pros y contras}
  \ProCon{Historia espacial muy fuerte\\Datos oficiales ESA\\Factible en 6 semanas\\Buen oral}{Casos peligrosos muy raros\\Ya hay resultados de competición\\Pocas imágenes ``bonitas''\\Hay que aportar ángulo propio}
\end{frame}

\begin{frame}{1 · Impacto (con los pies en el suelo)}
  \begin{itemize}
    \item Mejor aviso temprano $\rightarrow$ menos maniobras inútiles y más seguridad.
    \item No pretendemos sustituir el control de misión de la ESA: es un estudio académico con datos históricos.
  \end{itemize}
\end{frame}

\begin{frame}{1 · Veredicto}
  \Verdict{Muy recomendada (prioridad alta)}{Ideal si queréis que el proyecto se lea como ``espacio / ESA''. Alternativas cercanas: dinámica orbital (nº 4) o incendios por satélite (nº 5).}
\end{frame}

\subsection{Opción 2 --- PHBench}

\begin{frame}{2 · Portada --- ¿Qué startups llegan a Series A?}
  \OptionHeader{Recomendada nº 2}{Inversión / VC}{%
    Usamos señales públicas del día del lanzamiento en Product Hunt para estimar qué productos acaban levantando una ronda Series A.
  }
\end{frame}

\begin{frame}{2 · Imagina esto (sin jerga)}
  \begin{itemize}
    \item Cada día salen cientos de apps y productos nuevos.
    \item Casi ninguno consigue una gran ronda de inversión (Series A) en el año y medio siguiente.
    \item Product Hunt es como un ``escaparate'' público: votos, comentarios, equipo, descripción\ldots
    \item Un inversor no puede mirarlo todo: necesita una \textbf{lista ordenada} de ``esto merece 10 minutos''.
    \item Nosotros entrenamos modelos que intentan hacer ese ranking con datos pasados.
  \end{itemize}
\end{frame}

\begin{frame}{2 · Analogía cotidiana}
  \begin{center}
  \begin{tikzpicture}
    \node[box, text width=3.2cm] (a) {Muchas\\candidaturas};
    \node[box, text width=3.2cm, right=0.6cm of a] (b) {Pocas\\plazas};
    \node[boxdark, text width=3.2cm, right=0.6cm of b] (c) {Hay que\\priorizar};
    \draw[arr] (a)--(b)--(c);
  \end{tikzpicture}
  \end{center}
  \vspace{0.5em}
  Es el mismo tipo de problema que filtrar muchas solicitudes cuando solo unas pocas son ``éxito''.\\
  \textbf{No} es predecir si sube o baja una acción en bolsa.
\end{frame}

\begin{frame}{2 · ¿Qué haría aquí el machine learning?}
  \begin{itemize}
    \item Aprender patrones en lanzamientos antiguos (cuántos votos, qué temas, qué texto\ldots).
    \item Predecir la probabilidad de Series A --- sabiendo que el éxito es \textbf{rarísimo} ($\sim$1 de cada 130).
    \item Evaluar con métricas de ranking: ``de los 50 que marcamos como prometedores, ¿cuántos acertamos?''
    \item Cuidado crítico: no usar información del futuro (fugas temporales).
  \end{itemize}
\end{frame}

\begin{frame}{2 · Glosario rápido}
  \begin{description}[font=\bfseries\scriptsize,style=nextline,leftmargin=1.4cm]
    \item[Series A] Primera gran ronda institucional de una startup.
    \item[Product Hunt] Comunidad donde se lanzan productos digitales y se votan.
    \item[Desbalance] Casi todo es ``no Series A''; acertar el ``no'' es fácil y engañoso.
    \item[PHBench] Benchmark abierto (2026) con launches etiquetados.
  \end{description}
\end{frame}

\begin{frame}{2 · Pregunta de investigación}
  \begin{block}{Pregunta principal}
    Con señales del día del launch, ¿qué modelos ordenan mejor los candidatos a Series A en 18 meses, respetando el orden temporal de los datos?
  \end{block}
  \begin{block}{Pregunta secundaria}
    ¿El texto de la descripción aporta algo más allá de los números (votos, ranking del día\ldots)?
  \end{block}
\end{frame}

\begin{frame}{2 · Cómo lo haríamos, paso a paso}
  \PipelineFive{Orden\\temporal}{Variables\\+ texto}{Modelos\\clásicos}{Calibrar\\probabilidades}{Analizar\\errores}
\end{frame}

\begin{frame}{2 · Datos}
  \begin{itemize}
    \item $\sim$67\,000 lanzamientos (2019--2025), $\sim$61 variables documentadas.
    \item Etiquetas de Series A auditadas; hay clasificación pública de referencia.
    \item Acción inmediata: confirmar acceso al dataset (Hugging Face).
  \end{itemize}
\end{frame}

\begin{frame}{2 · Métodos del curso y más allá}
  \MethodsSplit{Logística regularizada · Árboles / boosting · Validación temporal · Métricas de ranking}{Representaciones de texto · Ensambles · Calibración · Análisis de qué variables importan}
\end{frame}

\begin{frame}{2 · Por qué encaja con el curso}
  \begin{columns}[T]
    \begin{column}{0.55\textwidth}
      \begin{itemize}
        \item Original frente a crédito/bolsa típicos.
        \item Obliga a rigor con eventos raros.
        \item Muy enseñable en entrevistas.
      \end{itemize}
    \end{column}
    \begin{column}{0.42\textwidth}
      \ScoreBar{Ambición}{9}{ieblue}
      \ScoreBar{Datos}{9}{accent}
      \ScoreBar{Factibilidad}{8}{good}
      \ScoreBar{Portfolio}{10}{good}
    \end{column}
  \end{columns}
\end{frame}

\begin{frame}{2 · Por qué encaja con nosotros}
  \FitDual{Le interesa el mundo VC y rechaza finanzas ``de siempre''. Buen portfolio.}{Finanzas + sistemas de recomendación/ranking. Quiere datos claros y comparación limpia.}{8}{10}
\end{frame}

\begin{frame}{2 · Pros y contras}
  \ProCon{Poco genérico\\Encaja con los dos\\Baselines listos\\Historia fuerte en entrevista}{Éxitos muy raros\\Acceso a datos por confirmar\\Fácil contaminar con futuro\\Menos ``deeptech físico''}
\end{frame}

\begin{frame}{2 · Impacto}
  Ayudar a priorizar atención en early-stage con datos públicos. No es consejo financiero: es un ranking empírico con error.
\end{frame}

\begin{frame}{2 · Veredicto}
  \Verdict{Muy recomendada (prioridad alta)}{Ideal si la identidad es ``inversión / portfolio''. No la combineéis como núcleo con VCBench (nº 7).}
\end{frame}

\subsection{Opción 3 --- SCAN}

\begin{frame}{3 · Portada --- ¿Los modelos ``entienden'' o solo memorizan?}
  \OptionHeader{Recomendada nº 3}{Fundamentos / lenguaje}{%
    Medimos si un modelo que ``conoce'' una instrucción suelta es capaz de recombinarla en frases nuevas --- el famoso fallo tipo ``can't jump''.
  }
\end{frame}

\begin{frame}{3 · Imagina esto (sin jerga)}
  \begin{itemize}
    \item A un niño le enseñas el verbo nuevo ``dax'' y entiende enseguida ``dax dos veces''.
    \item Muchos modelos de IA, en cambio, fallan en esa recombinación.
    \item No es un debate filosófico vacío: se puede \textbf{medir} con un examen diseñado para eso.
    \item SCAN es ese examen: órdenes simples (``salta'', ``anda'', ``dos veces'') convertidas en acciones.
  \end{itemize}
\end{frame}

\begin{frame}{3 · Analogía cotidiana}
  \begin{center}
  \begin{tikzpicture}
    \node[box, text width=2.6cm] (a) {Sabes\\``saltar''};
    \node[box, text width=2.6cm, right=0.45cm of a] (b) {Sabes\\``dos veces''};
    \node[boxdark, text width=3.0cm, right=0.45cm of b] (c) {¿Entiendes\\``salta 2 veces''?};
    \draw[arr] (a)--(b)--(c);
  \end{tikzpicture}
  \end{center}
  \vspace{0.45em}
  Si solo memorizaste ejemplos raros del examen, puedes aprobar el tipo A y suspender el tipo B.\\
  Eso es \textbf{generalización composicional}.
\end{frame}

\begin{frame}{3 · ¿Qué haría aquí el machine learning?}
  \begin{itemize}
    \item Entrenar modelos clásicos y redes en un trozo de SCAN.
    \item Probarlos en el split donde ``jump'' casi no aparece compuesto.
    \item Comparar: ¿quién falla y en qué tipo de órdenes?
    \item Opcional: repetir la idea en un dataset más lingüístico (COGS) para responder ``¿SCAN ya está resuelto?''.
  \end{itemize}
\end{frame}

\begin{frame}{3 · Glosario rápido}
  \begin{description}[font=\bfseries\scriptsize,style=nextline,leftmargin=1.5cm]
    \item[Composición] Combinar piezas conocidas en algo nuevo.
    \item[OOD] Fuera de lo visto en entrenamiento (el examen ``difícil'').
    \item[SCAN] Banco de pruebas de órdenes $\rightarrow$ acciones (Lake \& Baroni).
    \item[addprim\_jump] Split concreto donde ``jump'' solo se ve aislado en train.
  \end{description}
\end{frame}

\begin{frame}{3 · Pregunta de investigación}
  \begin{block}{Pregunta}
    Cuando el modelo ve ``jump'' solo en aislamiento, ¿qué familias de métodos fallan al usarlo dentro de órdenes nuevas, y qué señales del input predicen ese fallo?
  \end{block}
\end{frame}

\begin{frame}{3 · Cómo lo haríamos, paso a paso}
  \PipelineFive{Splits\\oficiales}{Modelos\\clásicos}{Redes\\secuencia}{Análisis\\de errores}{COGS\\opcional}
\end{frame}

\begin{frame}{3 · Datos}
  \begin{itemize}
    \item SCAN: pequeño, limpio, se corre en un portátil.
    \item Paper clásico (ICML 2018) = bibliografía fuerte.
    \item Defensa ante ``¿es realista?'': es el banco de pruebas estándar de composición, como un test de conducción controlado.
  \end{itemize}
\end{frame}

\begin{frame}{3 · Métodos del curso y más allá}
  \MethodsSplit{N-gramas / TF-IDF · Logística / SVM · Árboles · Validación OOD}{LSTM/GRU · Transformer pequeño · Análisis de errores · Modelo de lenguaje abierto solo como contraste}
\end{frame}

\begin{frame}{3 · Por qué encaja con el curso}
  \begin{columns}[T]
    \begin{column}{0.55\textwidth}
      \begin{itemize}
        \item Muy original frente a Kaggle típico.
        \item La comparación de métodos es el núcleo.
        \item Oral memorable si se cuenta bien.
      \end{itemize}
    \end{column}
    \begin{column}{0.42\textwidth}
      \ScoreBar{Ambición}{9}{ieblue}
      \ScoreBar{Factibilidad}{9}{good}
      \ScoreBar{Diferenciación}{10}{accent}
      \ScoreBar{Portfolio VC}{5}{bad}
    \end{column}
  \end{columns}
\end{frame}

\begin{frame}{3 · Por qué encaja con nosotros}
  \FitDual{Quiere profundidad tipo investigación y temas de abducción / límites de los LLM.}{Le gusta NLP y comparar modelos clásicos frente a redes con análisis claro.}{10}{8}
\end{frame}

\begin{frame}{3 · Pros y contras}
  \ProCon{Muy diferenciador\\Muy factible\\Oral fuerte\\Encaja con los dos}{Datos sintéticos\\Crítica ``ya está estudiado''\\Menos portfolio espacial/VC\\Conviene añadir COGS o diagnóstico}
\end{frame}

\begin{frame}{3 · Impacto}
  Aclarar cuándo un modelo parece inteligente pero no recombina reglas --- útil para diseñar evaluaciones honestas.
\end{frame}

\begin{frame}{3 · Veredicto}
  \Verdict{Muy recomendada (prioridad alta)}{Ideal si la identidad es ``fundamentos / composición''. Alternativa cercana: SCAR (nº 8).}
\end{frame}

\subsection{Opción 4 --- Orbital HNN / PINN}

\begin{frame}{4 · Portada --- Enseñar física a la red neuronal}
  \OptionHeader{Recomendada nº 4}{Física + ML}{%
    Entrenamos modelos que predicen trayectorias, pero obligándolos a respetar reglas físicas (energía, momento), y medimos si extrapolan mejor.
  }
\end{frame}

\begin{frame}{4 · Imagina esto (sin jerga)}
  \begin{itemize}
    \item Una red neuronal puede ``copiar'' una órbita en los datos de entrenamiento.
    \item Cuando le pides el futuro un poco más lejos, a veces \textbf{se desmadra}: la trayectoria deja de parecer física.
    \item Un ingeniero sabe que ciertas cantidades no deberían ``desaparecer'' (energía, momento angular).
    \item La idea: meter esas reglas dentro del aprendizaje, no solo mirar el error numérico.
  \end{itemize}
\end{frame}

\begin{frame}{4 · Analogía cotidiana}
  \begin{center}
  \begin{tikzpicture}
    \node[box, text width=3.4cm] (a) {Aprender a tirar\\una pelota};
    \node[box, text width=3.4cm, right=0.5cm of a] (b) {Sin gravedad\\en la cabeza};
    \node[boxdark, text width=3.4cm, right=0.5cm of b] (c) {Con gravedad\\como regla};
    \draw[arr] (a)--(b)--(c);
  \end{tikzpicture}
  \end{center}
  \vspace{0.4em}
  Con la regla, suele hacer falta \textbf{menos datos} y el modelo aguanta mejor situaciones nuevas.\\
  Condición: debe haber un ancla en datos reales (no solo un ejercicio de libro).
\end{frame}

\begin{frame}{4 · ¿Qué haría aquí el machine learning?}
  \begin{itemize}
    \item Comparar: regresión clásica / red normal / red con física (PINN o red hamiltoniana).
    \item Medir dos cosas: ¿acota bien la posición? ¿respeta la física (poca deriva de energía)?
    \item Hacer ablación: apagar la física y ver qué se rompe.
  \end{itemize}
\end{frame}

\begin{frame}{4 · Glosario rápido}
  \begin{description}[font=\bfseries\scriptsize,style=nextline,leftmargin=1.6cm]
    \item[PINN] Red que penaliza incumplir una ecuación física.
    \item[HNN] Red que aprende una energía y avanza de forma más ``física''.
    \item[Extrapolación] Predecir fuera del intervalo visto.
    \item[Ablación] Quitar una pieza del modelo para ver si importaba.
  \end{description}
\end{frame}

\begin{frame}{4 · Pregunta de investigación}
  \begin{block}{Pregunta}
    ¿Un modelo con reglas físicas reduce el error al extrapolar y conserva mejor energía/momento que una red sin esas reglas, con pocos datos ruidosos?
  \end{block}
\end{frame}

\begin{frame}{4 · Cómo lo haríamos, paso a paso}
  \PipelineFive{Elegir\\sistema}{Modelos\\base}{PINN\\o HNN}{Medir\\error+física}{Ablación\\del prior}
\end{frame}

\begin{frame}{4 · Datos posibles}
  \begin{itemize}
    \item Simulación controlada + ruido (para entender).
    \item Ancla real: catálogos orbitales, velocidades radiales de exoplanetas, o un péndulo medido.
    \item En la proposal hay que nombrar ya cuál será el ancla.
  \end{itemize}
\end{frame}

\begin{frame}{4 · Métodos del curso y más allá}
  \MethodsSplit{Regresión · Redes del curso · Regularización · Validación}{PINN · Redes hamiltonianas · Residuo con GP · Latentes con significado físico}
\end{frame}

\begin{frame}{4 · Por qué encaja con el curso}
  \begin{columns}[T]
    \begin{column}{0.55\textwidth}
      \begin{itemize}
        \item Ambición metodológica (no solo accuracy).
        \item Muy alineado con el estilo del profesor (mecanismo + comparación).
        \item Oral exigente pero brillante si se domina.
      \end{itemize}
    \end{column}
    \begin{column}{0.42\textwidth}
      \ScoreBar{Ambición}{9}{ieblue}
      \ScoreBar{Alineación profe}{10}{good}
      \ScoreBar{Factibilidad}{7}{accent}
      \ScoreBar{Afinidad Mario}{6}{bad}
    \end{column}
  \end{columns}
\end{frame}

\begin{frame}{4 · Por qué encaja con nosotros}
  \FitDual{Física al máximo. Puede proponer la capa creativa (ruido, disipación, incertidumbre).}{Puede reproducir papers y llevar la comparación rigurosa. Menos ``producto de negocio''.}{10}{6}
\end{frame}

\begin{frame}{4 · Pros y contras}
  \ProCon{Muy alineado con el profe\\Pedro muy motivado\\Ablaciones potentes\\Historia científica limpia}{Riesgo de parecer demo de libro\\Mario menos enganchado\\Hay que fijar dato real\\Menos portfolio VC}
\end{frame}

\begin{frame}{4 · Impacto}
  Mejores predicciones de dinámica con menos datos --- idea transferable a espacio, control y ciencia computacional.
\end{frame}

\begin{frame}{4 · Veredicto}
  \Verdict{Recomendada (con condición)}{Solo si en la proposal ya hay dato real + baselines. Ideal para identidad ``física / método del profe sin bio''.}
\end{frame}

\subsection{Opción 5 --- Wildfire}

\begin{frame}{5 · Portada --- Ver incendios desde el satélite (aunque haya nubes)}
  \OptionHeader{Recomendada nº 5}{Clima / satélites}{%
    Combinamos fotos ópticas y radar de satélites europeos para detectar o mapear fuego, y medimos cuánto aporta cada sensor.
  }
\end{frame}

\begin{frame}{5 · Imagina esto (sin jerga)}
  \begin{itemize}
    \item Desde el espacio se ve la Tierra en ``fotos'' (luz visible e infrarroja).
    \item Si hay nubes, esas fotos no sirven: el fuego queda tapado.
    \item Hay otro tipo de satélite que usa \textbf{radar}: no ``ve'' como una cámara, pero atraviesa nubes.
    \item Pregunta humana: ¿juntar cámara + radar ayuda de verdad a no perder incendios?
  \end{itemize}
\end{frame}

\begin{frame}{5 · Analogía cotidiana}
  \begin{center}
  \begin{tikzpicture}
    \node[box, text width=3.3cm] (a) {Cámara\\(día despejado)};
    \node[box, text width=3.3cm, right=0.45cm of a] (b) {Radar\\(niebla/nubes)};
    \node[boxdark, text width=3.3cm, right=0.45cm of b] (c) {Decisión\\mejor informada};
    \draw[arr] (a)--(b)--(c);
  \end{tikzpicture}
  \end{center}
  \vspace{0.4em}
  Como conducir con faros normales \emph{y} sensores de aparcamiento: cada uno falla en sitios distintos.
\end{frame}

\begin{frame}{5 · ¿Qué haría aquí el machine learning?}
  \begin{itemize}
    \item Calcular índices de vegetación/quemado a partir de bandas espectrales.
    \item Añadir señales de radar.
    \item Entrenar clasificadores (árboles, redes) para ``quemado / no quemado'' o fuego activo.
    \item Hacer ablaciones: quitar el radar y medir cuánto empeora, sobre todo con nubes.
  \end{itemize}
\end{frame}

\begin{frame}{5 · Glosario rápido}
  \begin{description}[font=\bfseries\scriptsize,style=nextline,leftmargin=1.8cm]
    \item[Sentinel-2] Satélite óptico europeo (``cámara'' multiespectral).
    \item[Sentinel-1] Satélite radar europeo (funciona con nubes).
    \item[Burnt area] Zona ya quemada.
    \item[Fusión] Usar varias fuentes a la vez, de forma transparente.
  \end{description}
\end{frame}

\begin{frame}{5 · Pregunta de investigación}
  \begin{block}{Pregunta}
    ¿Cuánto mejora la detección al fusionar óptico + radar frente a usar solo uno, y qué variables importan más bajo nubes?
  \end{block}
\end{frame}

\begin{frame}{5 · Cómo lo haríamos, paso a paso}
  \PipelineFive{Dataset\\listo}{Variables\\S2/S1}{Árboles /\\boosting}{Red\\opcional}{Ablación\\sensores}
  \vspace{0.3em}
  Empezar con datasets empaquetados (p.\ ej.\ Sen2Fire), no con descargas masivas día 1.
\end{frame}

\begin{frame}{5 · Datos}
  \begin{itemize}
    \item Sen2Fire y similares: recortes ya preparados.
    \item Alternativas: productos europeos de área quemada.
    \item Riesgo: si escaláis a cubos enormes, el almacenamiento se come el tiempo.
  \end{itemize}
\end{frame}

\begin{frame}{5 · Métodos del curso y más allá}
  \MethodsSplit{Árboles / boosting · Logística · Ingeniería de variables · Validación espacial}{CNN / U-Net · Fusión de sensores · Desbalance de clases · Mapas de incertidumbre}
\end{frame}

\begin{frame}{5 · Por qué encaja con el curso}
  \begin{columns}[T]
    \begin{column}{0.55\textwidth}
      \begin{itemize}
        \item Multi-fuente (valorado en el enunciado).
        \item Impacto fácil de explicar.
        \item Figuras muy útiles para Communication Skills.
      \end{itemize}
    \end{column}
    \begin{column}{0.42\textwidth}
      \ScoreBar{Ambición}{8}{ieblue}
      \ScoreBar{Impacto}{9}{good}
      \ScoreBar{Factibilidad}{6}{accent}
      \ScoreBar{Visuales}{9}{good}
    \end{column}
  \end{columns}
\end{frame}

\begin{frame}{5 · Por qué encaja con nosotros}
  \FitDual{Satélites y visión le motivan; busca impacto real.}{Clima le interesa; puede llevar comparación de modelos y métricas.}{8}{7}
\end{frame}

\begin{frame}{5 · Pros y contras}
  \ProCon{Impacto claro\\Historia multi-sensor\\Buenas figuras\\Encaje razonable de ambos}{El pipeline satelital come tiempo\\Menos ``misión espacial'' que ESA\\Etiquetas ruidosas\\Cuidado con el tamaño de datos}
\end{frame}

\begin{frame}{5 · Impacto}
  Mejor mapa de fuego en mal tiempo --- con humildad: detectar no es lo mismo que explicar causas climáticas.
\end{frame}

\begin{frame}{5 · Veredicto}
  \Verdict{Recomendada}{Buena opción ``clima/satélite''. Si el volumen asusta: OSCD (nº 9) solo con un giro propio.}
\end{frame}


\section{Condicionales}
\subsection{Opción 6 --- Matbench}

\begin{frame}{6 · Portada --- ¿Qué tan ``conductor'' es un material?}
  \OptionHeader{Condicional nº 6}{Materiales}{%
    Predecimos una propiedad clave de materiales (el band gap) con un banco de datos estándar, midiendo error e incertidumbre.
  }
\end{frame}

\begin{frame}{6 · Imagina esto (sin jerga)}
  \begin{itemize}
    \item Para inventar una batería, un chip o un panel, hay que probar muchos materiales.
    \item Probarlos en laboratorio o simularlos átomo a átomo es \textbf{caro y lento}.
    \item El \emph{band gap} resume, a grandes rasgos, si algo se comporta como metal, semiconductor o aislante.
    \item Si un modelo estima esa propiedad bien --- y dice cuándo duda --- se pueden priorizar candidatos.
  \end{itemize}
\end{frame}

\begin{frame}{6 · Analogía cotidiana}
  \begin{center}
  \begin{tikzpicture}
    \node[box, text width=3.5cm] (a) {Receta\\(composición)};
    \node[box, text width=3.5cm, right=0.5cm of a] (b) {Modelo\\que estima};
    \node[boxdark, text width=3.5cm, right=0.5cm of b] (c) {``Nota'' +\\``cómo de seguro''};
    \draw[arr] (a)--(b)--(c);
  \end{tikzpicture}
  \end{center}
  \vspace{0.35em}
  Como estimar calorías de un plato por ingredientes: útil si también marcas la incertidumbre.
\end{frame}

\begin{frame}{6 · ¿Qué haría aquí el machine learning?}
  \begin{itemize}
    \item Convertir la fórmula química (o estructura) en variables numéricas.
    \item Comparar regresiones clásicas, árboles y redes.
    \item Reportar no solo el error medio: también \textbf{cuándo} el modelo está perdido.
  \end{itemize}
\end{frame}

\begin{frame}{6 · Glosario rápido}
  \begin{description}[font=\bfseries\scriptsize,style=nextline,leftmargin=1.8cm]
    \item[Band gap] Escala de energía relacionada con el comportamiento electrónico.
    \item[Matbench] Banco de tareas estándar de ML para materiales.
    \item[UQ] Cuantificar incertidumbre (``no solo el número, también la duda'').
  \end{description}
\end{frame}

\begin{frame}{6 · Pregunta de investigación}
  \begin{block}{Pregunta}
    ¿Qué combinación de representaciones + modelos minimiza el error de band gap con incertidumbre calibrada, y en qué familias químicas falla más?
  \end{block}
\end{frame}

\begin{frame}{6 · Cómo lo haríamos, paso a paso}
  \PipelineFive{Elegir\\tarea}{Crear\\variables}{Árboles /\\boosting}{Red +\\incertidumbre}{Errores por\\química}
\end{frame}

\begin{frame}{6 · Datos}
  \begin{itemize}
    \item Matbench: tareas limpias y citables.
    \item Empezar por la versión experimental más pequeña ($\sim$4\,600); escalar si sobra tiempo.
  \end{itemize}
\end{frame}

\begin{frame}{6 · Métodos del curso y más allá}
  \MethodsSplit{Regresión · Árboles / boosting · Kernels · Validación anidada}{Multi-fidelidad · Conformal / GP · Análisis por familia química}
\end{frame}

\begin{frame}{6 · Por qué encaja con el curso}
  \begin{columns}[T]
    \begin{column}{0.55\textwidth}
      \begin{itemize}
        \item Datos impecables.
        \item Muy articulable con el temario.
        \item La incertidumbre aporta ambición sobria.
      \end{itemize}
    \end{column}
    \begin{column}{0.42\textwidth}
      \ScoreBar{Ambición}{7}{accent}
      \ScoreBar{Datos}{10}{good}
      \ScoreBar{Factibilidad}{9}{good}
      \ScoreBar{Wow Pedro}{6}{bad}
    \end{column}
  \end{columns}
\end{frame}

\begin{frame}{6 · Por qué encaja con nosotros}
  \FitDual{Ciencia dura OK, pero menos emoción ``espacial''. Hace falta el giro de incertidumbre o se queda corto.}{Buen encaje con trabajo limpio, comparación y resultados claros.}{6}{7}
\end{frame}

\begin{frame}{6 · Pros y contras}
  \ProCon{Muy factible\\Datos excelentes\\Fácil de defender\\UQ seria}{Menos impacto ``wow''\\Riesgo de parecer solo un ranking\\Mario no viene de materials\\Pedro puede pedirlo más ambicioso}
\end{frame}

\begin{frame}{6 · Impacto}
  Acelerar el filtrado de materiales candidatos --- la aportación es método + duda calibrada, no un material nuevo.
\end{frame}

\begin{frame}{6 · Veredicto}
  \Verdict{Condicional --- excelente plan B}{Elegidla si priorizáis ejecución limpia y nota sólida. Subid ambición con incertidumbre.}
\end{frame}

\subsection{Opción 7 --- VCBench}

\begin{frame}{7 · Portada --- ¿El perfil del founder anticipa un gran éxito?}
  \OptionHeader{Condicional nº 7}{Inversión / founders}{%
    Con perfiles anonimizados de fundadores (antes de fundar), intentamos estimar la probabilidad de un éxito empresarial muy grande.
  }
\end{frame}

\begin{frame}{7 · Imagina esto (sin jerga)}
  \begin{itemize}
    \item Un inversor early-stage ve currículums y trayectorias, no el futuro.
    \item Algunas personas tienen historial (estudios, empresas previas, sector\ldots).
    \item VCBench pregunta: con \textbf{solo} lo que se sabía antes de fundar, ¿se puede anticipar un éxito enorme?
    \item ``Éxito enorme'' aquí está definido de forma dura (muy alto funding o exit).
  \end{itemize}
\end{frame}

\begin{frame}{7 · Analogía cotidiana}
  \begin{center}
  \begin{tikzpicture}
    \node[box, text width=3.4cm] (a) {Historial\\antes del partido};
    \node[box, text width=3.4cm, right=0.5cm of a] (b) {Modelo\\de probabilidad};
    \node[boxdark, text width=3.4cm, right=0.5cm of b] (c) {¿Alto\\rendimiento?};
    \draw[arr] (a)--(b)--(c);
  \end{tikzpicture}
  \end{center}
  \vspace{0.35em}
  Importante: no mirar el marcador final para ``predecir'' el marcador final (fuga de información).
\end{frame}

\begin{frame}{7 · ¿Qué haría aquí el machine learning?}
  \begin{itemize}
    \item Mezclar variables estructuradas + texto del perfil.
    \item Comparar modelos y mirar precisión en los positivos (porque un falso ``sí'' es caro).
    \item Analizar sesgos: ¿el modelo premia proxies injustos?
  \end{itemize}
\end{frame}

\begin{frame}{7 · Glosario rápido}
  \begin{description}[font=\bfseries\scriptsize,style=nextline,leftmargin=1.6cm]
    \item[Founder] Persona fundadora de la startup.
    \item[Major success] Etiqueta del benchmark (umbral muy alto de funding/exit).
    \item[F0.5] Métrica que castiga más los falsos positivos que el recall puro.
  \end{description}
\end{frame}

\begin{frame}{7 · Pregunta de investigación}
  \begin{block}{Pregunta}
    ¿Qué modelos maximizan precisión / F0.5 en éxitos mayores, y qué partes del historial importan sin usar información futura?
  \end{block}
\end{frame}

\begin{frame}{7 · Cómo lo haríamos, paso a paso}
  \PipelineFive{Kit de\\datos}{Tabla\\+ texto}{Comparar\\modelos}{Precisión\\/ F0.5}{Sesgos y\\errores}
\end{frame}

\begin{frame}{7 · Datos}
  \begin{itemize}
    \item $\sim$4\,500 perfiles públicos de entrenamiento.
    \item $\sim$9\% positivos (más manejable que PHBench).
    \item Hay kit de inicio en GitHub.
  \end{itemize}
\end{frame}

\begin{frame}{7 · Métodos del curso y más allá}
  \MethodsSplit{Logística / árboles · Regularización · Desbalance · Validación}{Embeddings de texto · Modelo híbrido · Calibración · Revisión de sesgos}
\end{frame}

\begin{frame}{7 · Por qué encaja con el curso}
  \begin{columns}[T]
    \begin{column}{0.55\textwidth}
      \begin{itemize}
        \item Original frente a crédito genérico.
        \item Buena data + comparación.
        \item Ética y limitaciones naturales.
      \end{itemize}
    \end{column}
    \begin{column}{0.42\textwidth}
      \ScoreBar{Ambición}{8}{ieblue}
      \ScoreBar{Datos}{8}{accent}
      \ScoreBar{Factibilidad}{8}{good}
      \ScoreBar{Solapa nº 2}{9}{bad}
    \end{column}
  \end{columns}
\end{frame}

\begin{frame}{7 · Por qué encaja con nosotros}
  \FitDual{VC innovador OK; útil si PHBench no es accesible a tiempo.}{Finanzas + texto. Quiere análisis limpio.}{8}{10}
\end{frame}

\begin{frame}{7 · Pros y contras}
  \ProCon{CSV usable enseguida\\Menos extremo que PHBench\\Historia VC clara}{Solapa con la opción 2\\Menos datos\\Test privado\\Hay que hablar de sesgos}
\end{frame}

\begin{frame}{7 · Impacto}
  Entender límites de ``predecir éxito'' de personas/empresas: ranking con humildad, no oráculo.
\end{frame}

\begin{frame}{7 · Veredicto}
  \Verdict{Condicional --- respaldo de PHBench}{Si falla el acceso a PHBench, pivotad aquí. No dos núcleos VC a la vez.}
\end{frame}

\subsection{Opción 8 --- SCAR}

\begin{frame}{8 · Portada --- ¿La IA entiende analogías de verdad?}
  \OptionHeader{Condicional nº 8}{Lenguaje / abducción}{%
    Medimos si un modelo no solo elige una analogía ``parecida'', sino que alinea correctamente las piezas de dos sistemas (mapa estructural).
  }
\end{frame}

\begin{frame}{8 · Imagina esto (sin jerga)}
  \begin{itemize}
    \item Decimos: ``el átomo es como un sistema solar''.
    \item No basta con decir ``sí, hay analogía''.
    \item Hay que mapear: Sol$\leftrightarrow$núcleo, planetas$\leftrightarrow$electrones,\ldots
    \item SCAR es un examen de ese mapeo en analogías científicas.
    \item Muchos modelos fallan ahí aunque acierten preguntas fáciles de parecido.
  \end{itemize}
\end{frame}

\begin{frame}{8 · Analogía cotidiana}
  \begin{center}
  \begin{tikzpicture}
    \node[box, text width=3.3cm] (a) {Equipo de\\fútbol};
    \node[box, text width=3.3cm, right=0.5cm of a] (b) {Orquesta};
    \node[boxdark, text width=3.5cm, below=0.55cm of $(a)!0.5!(b)$] (c) {entrenador$\leftrightarrow$director\\jugadores$\leftrightarrow$músicos};
    \draw[arr] (a)|-(c); \draw[arr] (b)|-(c);
  \end{tikzpicture}
  \end{center}
  \vspace{0.35em}
  Entender la analogía es construir el diccionario entre piezas, no solo sentir que ``se parecen''.
\end{frame}

\begin{frame}{8 · ¿Qué haría aquí el machine learning?}
  \begin{itemize}
    \item Representar conceptos con embeddings.
    \item Proponer emparejamientos (matching) y puntuarlos.
    \item Comparar con reglas simples de solapamiento de palabras.
    \item Usar un LLM grande solo como contraste, no como autor del report.
  \end{itemize}
\end{frame}

\begin{frame}{8 · Glosario rápido}
  \begin{description}[font=\bfseries\scriptsize,style=nextline,leftmargin=1.8cm]
    \item[Abducción] Inferir la mejor explicación/estructura a partir de indicios.
    \item[Mapping] Correspondencia pieza a pieza entre dos sistemas.
    \item[SCAR] Benchmark de analogías científicas con estructura.
  \end{description}
\end{frame}

\begin{frame}{8 · Pregunta de investigación}
  \begin{block}{Pregunta}
    ¿Qué representaciones y algoritmos de emparejamiento recuperan mejor los mappings de SCAR, y en qué dominios fallan frente a un LLM en few-shot?
  \end{block}
\end{frame}

\begin{frame}{8 · Cómo lo haríamos, paso a paso}
  \PipelineFive{Cargar\\SCAR}{Baselines\\léxicos}{Embeddings\\+ matching}{Contraste\\LLM}{Errores por\\dominio}
\end{frame}

\begin{frame}{8 · Datos}
  \begin{itemize}
    \item 400 analogías, 13 dominios, 1600 mappings.
    \item Tamaño ideal para un proyecto de curso.
  \end{itemize}
\end{frame}

\begin{frame}{8 · Métodos del curso y más allá}
  \MethodsSplit{Kernels / SVM · Árboles · PCA / similitud · Clasificación}{Embeddings de frases · Matching · Transformer pequeño · LLM solo como contraste}
\end{frame}

\begin{frame}{8 · Por qué encaja con el curso}
  \begin{columns}[T]
    \begin{column}{0.55\textwidth}
      \begin{itemize}
        \item Muy original.
        \item Bibliografía rica.
        \item Oral intelectual fuerte.
      \end{itemize}
    \end{column}
    \begin{column}{0.42\textwidth}
      \ScoreBar{Ambición}{9}{ieblue}
      \ScoreBar{Profundidad NLP}{9}{accent}
      \ScoreBar{Factibilidad}{6}{accent}
      \ScoreBar{Portfolio}{6}{bad}
    \end{column}
  \end{columns}
\end{frame}

\begin{frame}{8 · Por qué encaja con nosotros}
  \FitDual{Encaja con abducción / límites del lenguaje si el problema es ambicioso.}{NLP alto. Puede llevar métodos y análisis.}{9}{9}
\end{frame}

\begin{frame}{8 · Pros y contras}
  \ProCon{Muy temático con abducción\\Diferenciador\\Dataset manejable}{Más pesado que SCAN\\Menos portfolio ``producto''\\Evaluar matching es delicado\\Solapa con la opción 3}
\end{frame}

\begin{frame}{8 · Impacto}
  Mejores pruebas de razonamiento analógico --- para evaluar modelos con más honestidad.
\end{frame}

\begin{frame}{8 · Veredicto}
  \Verdict{Condicional --- respaldo de SCAN}{Si queréis abducción más ``paper de NLP''. No dos núcleos 3 y 8 a la vez.}
\end{frame}

\subsection{Opción 9 --- OSCD}

\begin{frame}{9 · Portada --- ¿Qué cambió en la ciudad entre dos fotos satélite?}
  \OptionHeader{Condicional nº 9}{Satélite / ciudades}{%
    Detectamos cambios urbanos comparando dos imágenes Sentinel-2 del mismo lugar en fechas distintas.
  }
\end{frame}

\begin{frame}{9 · Imagina esto (sin jerga)}
  \begin{itemize}
    \item Google Earth, pero con propósito científico: dos fotos del mismo barrio, años distintos.
    \item Queremos un mapa que diga: aquí había campo y ahora hay edificios; aquí abrieron una carretera\ldots
    \item Eso sirve a urbanismo, seguros, ONGs, gobiernos.
    \item El peligro académico: el tema está muy visto --- sin un giro propio parece proyecto genérico.
  \end{itemize}
\end{frame}

\begin{frame}{9 · Analogía cotidiana}
  \begin{center}
  \begin{tikzpicture}
    \node[box, text width=2.8cm, minimum height=1.5cm] (a) {Foto\\antes};
    \node[box, text width=2.8cm, minimum height=1.5cm, right=0.55cm of a] (b) {Foto\\después};
    \node[boxdark, text width=3.0cm, right=0.55cm of b] (c) {Mapa de\\``qué cambió''};
    \draw[arr] (a)--(b)--(c);
  \end{tikzpicture}
  \end{center}
  \vspace{0.4em}
  Como el juego de ``encuentra las 7 diferencias'', pero automático y sobre ciudades enteras.
\end{frame}

\begin{frame}{9 · ¿Qué haría aquí el machine learning?}
  \begin{itemize}
    \item Baseline sencillo: diferencias de índices de vegetación + bosque aleatorio.
    \item Modelo más flexible: red que mira el par de imágenes a la vez.
    \item Giro obligatorio: importancia de bandas, incertidumbre, o transfer entre ciudades.
  \end{itemize}
\end{frame}

\begin{frame}{9 · Glosario rápido}
  \begin{description}[font=\bfseries\scriptsize,style=nextline,leftmargin=1.5cm]
    \item[Change detection] Marcar píxeles que cambiaron entre dos fechas.
    \item[OSCD] Dataset pequeño y clásico de cambio urbano con Sentinel-2.
    \item[Siamesa] Red que procesa las dos imágenes de forma emparejada.
  \end{description}
\end{frame}

\begin{frame}{9 · Pregunta de investigación}
  \begin{block}{Pregunta}
    ¿Una red emparejada supera a índices + bosque aleatorio en OSCD, y qué bandas importan más si exigimos incertidumbre o transferencia entre ciudades?
  \end{block}
\end{frame}

\begin{frame}{9 · Cómo lo haríamos, paso a paso}
  \PipelineFive{Bajar\\OSCD}{Baseline\\índices}{Red\\emparejada}{Ablación\\bandas}{Giro\\UQ/transfer}
\end{frame}

\begin{frame}{9 · Datos}
  \begin{itemize}
    \item 24 pares de imágenes; etiquetas en 14.
    \item Muy manejable, pero poca muestra $\Rightarrow$ cuidado con conclusiones.
  \end{itemize}
\end{frame}

\begin{frame}{9 · Métodos del curso y más allá}
  \MethodsSplit{Árboles · Ingeniería de variables · Clasificación · Validación}{CNN emparejada · Ablación de bandas · Incertidumbre · Transfer entre ciudades}
\end{frame}

\begin{frame}{9 · Por qué encaja con el curso}
  \begin{columns}[T]
    \begin{column}{0.55\textwidth}
      \begin{itemize}
        \item Datos reales de satélite.
        \item Camino claro curso $\rightarrow$ PyTorch.
        \item Figuras fuertes.
      \end{itemize}
    \end{column}
    \begin{column}{0.42\textwidth}
      \ScoreBar{Ambición}{6}{accent}
      \ScoreBar{Factibilidad}{9}{good}
      \ScoreBar{Riesgo genérico}{8}{bad}
      \ScoreBar{Visuales}{8}{good}
    \end{column}
  \end{columns}
\end{frame}

\begin{frame}{9 · Por qué encaja con nosotros}
  \FitDual{Satélites OK, pero necesita un giro o lo verá genérico.}{Visión media. Comparación OK. Menos pasión que finanzas.}{7}{6}
\end{frame}

\begin{frame}{9 · Pros y contras}
  \ProCon{Muy factible\\Dataset clásico\\Buenas visualizaciones\\Fácil de explicar}{Tema saturado\\Pocos pares\\Riesgo de verse genérico\\Menos diferenciación que ESA/incendios}
\end{frame}

\begin{frame}{9 · Impacto}
  Monitorizar expansión urbana --- el valor real depende del giro (bandas, incertidumbre, transferencia).
\end{frame}

\begin{frame}{9 · Veredicto}
  \Verdict{Condicional --- solo con giro propio}{Respaldo EO si incendios es demasiado pesado. Sin giro: mejor no.}
\end{frame}

\subsection{Opción 10 --- OXE subset}

\begin{frame}{10 · Portada --- Reutilizar experiencias de robots}
  \OptionHeader{Condicional nº 10}{Robótica}{%
    En un trozo pequeño de un gran dataset de robots, aprendemos a recomendar demostraciones útiles (o a predecir éxito/fallo) --- sin entrenar un modelo gigante.
  }
\end{frame}

\begin{frame}{10 · Imagina esto (sin jerga)}
  \begin{itemize}
    \item Hay muchos vídeos/demostraciones de brazos robot haciendo tareas (coger, colocar,\ldots).
    \item Cuando llega una tarea nueva, ¿qué demostraciones antiguas se parecen más?
    \item Eso es un \textbf{sistema de recomendación} sobre experiencias robot.
    \item El dataset completo es enorme: si intentáis usarlo entero, el proyecto se come en infraestructura.
  \end{itemize}
\end{frame}

\begin{frame}{10 · Analogía cotidiana}
  \begin{center}
  \begin{tikzpicture}
    \node[box, text width=3.3cm] (a) {Nueva receta\\que quiero hacer};
    \node[box, text width=3.3cm, right=0.5cm of a] (b) {Biblioteca de\\vídeos de cocina};
    \node[boxdark, text width=3.3cm, right=0.5cm of b] (c) {Top vídeos\\más útiles};
    \draw[arr] (a)--(b)--(c);
  \end{tikzpicture}
  \end{center}
  \vspace{0.35em}
  No reentrenáis YouTube entero: elegís un canal (subset) y medís si el ranking tiene sentido.
\end{frame}

\begin{frame}{10 · ¿Qué haría aquí el machine learning?}
  \begin{itemize}
    \item Congelar 1--2 subsets concretos (nombres en la proposal).
    \item Extraer descripciones numéricas de cada demostración.
    \item Comparar ranking clásico vs embeddings aprendidos.
    \item Métrica simple: de las K recomendaciones, ¿cuántas eran realmente útiles?
  \end{itemize}
\end{frame}

\begin{frame}{10 · Glosario rápido}
  \begin{description}[font=\bfseries\scriptsize,style=nextline,leftmargin=1.6cm]
    \item[OXE] Open X-Embodiment: colección enorme de demos de muchos robots.
    \item[Subset] Solo una parte pequeña y nombrada del dataset.
    \item[Recall@K] Entre las K primeras recomendaciones, cuántas eran correctas.
  \end{description}
\end{frame}

\begin{frame}{10 · Pregunta de investigación}
  \begin{block}{Pregunta}
    En subsets concretos, ¿representaciones que mezclan lenguaje, estado del robot y acciones recuperan mejores demostraciones que solo visión o variables simples?
  \end{block}
\end{frame}

\begin{frame}{10 · Cómo lo haríamos --- con freno de emergencia}
  \PipelineFive{Fijar\\subsets}{Extraer\\variables}{Baselines\\ranking}{Embeddings\\ligeros}{Recall@K}
  \vspace{0.35em}
  \begin{block}{Regla de parada}
    Si en 3 días no hay exploración usable por culpa de la infra, se cambia a un dataset de robótica más simple (ver brief largo: TrajAir, ALFA, etc.).
  \end{block}
\end{frame}

\begin{frame}{10 · Datos}
  \begin{itemize}
    \item Nunca el corpus completo (terabytes).
    \item Candidatos pequeños: Bridge, TOTO, etc.
    \item Los identificadores del subset deben estar en la proposal.
  \end{itemize}
\end{frame}

\begin{frame}{10 · Métodos del curso y más allá}
  \MethodsSplit{PCA / clustering · Árboles / logística · Autoencoders · Métricas de ranking}{Codificadores NN · VAE · Fusión multimodal · Recomendador de skills}
\end{frame}

\begin{frame}{10 · Por qué encaja con el curso}
  \begin{columns}[T]
    \begin{column}{0.55\textwidth}
      \begin{itemize}
        \item Ambición visible (robótica moderna).
        \item Recommenders = punto fuerte de Mario.
        \item Solo funciona si el alcance se respeta.
      \end{itemize}
    \end{column}
    \begin{column}{0.42\textwidth}
      \ScoreBar{Ambición}{8}{ieblue}
      \ScoreBar{Encaje equipo}{9}{good}
      \ScoreBar{Factibilidad}{5}{bad}
      \ScoreBar{Riesgo alcance}{9}{bad}
    \end{column}
  \end{columns}
\end{frame}

\begin{frame}{10 · Por qué encaja con nosotros}
  \FitDual{Robótica al máximo; historia cercana a labs de IE sin necesitar el lab.}{Recommenders + robótica. Necesita métricas y alcance férreos.}{9}{8}
\end{frame}

\begin{frame}{10 · Pros y contras}
  \ProCon{Mucho atractivo robótica\\Encaja con los dos\\Historia IE/CyPhy\\Recommenders reales}{Infra peligrosa\\El alcance se infla fácil\\Reproducir es frágil\\Oral difícil si el pipeline es caos}
\end{frame}

\begin{frame}{10 · Impacto}
  Reutilizar experiencia entre robots/tareas --- la idea grande de la robótica ``foundation'', en versión realista de curso.
\end{frame}

\begin{frame}{10 · Veredicto}
  \Verdict{Condicional --- solo con alcance congelado el día 1}{Si queréis robótica con menos riesgo, mirad alternativas del brief (TrajAir / ALFA / verificación + recomendación).}
\end{frame}


\section{Decisión}
\begin{frame}{Mapa de identidades: ¿qué historia queréis contar?}
  \begin{center}
  \begin{tikzpicture}
    \node[boxdark, text width=3.4cm] (a) {Espacio / ESA\\{\scriptsize $\to$ 1 (res. 4/5)}};
    \node[boxdark, text width=3.4cm, right=0.4cm of a] (b) {Inversión / VC\\{\scriptsize $\to$ 2 (res. 7)}};
    \node[boxdark, text width=3.4cm, right=0.4cm of b] (c) {Fundamentos IA\\{\scriptsize $\to$ 3 (res. 8)}};
    \node[box, text width=3.4cm, below=0.4cm of a] (d) {Materiales sólidos\\{\scriptsize $\to$ 6}};
    \node[box, text width=3.4cm, below=0.4cm of b] (e) {Clima / satélites\\{\scriptsize $\to$ 5 (res. 9)}};
    \node[box, text width=3.4cm, below=0.4cm of c] (f) {Robótica\\{\scriptsize $\to$ 10 solo si alcance fijo}};
  \end{tikzpicture}
  \end{center}
\end{frame}

\begin{frame}{Scorecard para rellenar juntos (1--5)}
  \tiny
  \renewcommand{\arraystretch}{1.25}
  \begin{tabular}{@{}clcccc@{}}
    \toprule
    \# & Opción & Afinidad Pedro & Afinidad Mario & Factib. 6 sem. & Curso \\
    \midrule
    1 & Satélites ESA & & & & \\
    2 & Series A / PHBench & & & & \\
    3 & Composición / SCAN & & & & \\
    4 & Física en la red & & & & \\
    5 & Incendios satélite & & & & \\
    6 & Matbench & & & & \\
    7 & VCBench founders & & & & \\
    8 & Analogías SCAR & & & & \\
    9 & Cambio urbano OSCD & & & & \\
    10 & Robots OXE (subset) & & & & \\
    \bottomrule
  \end{tabular}
  \vspace{0.45em}
  \normalsize Descartad cualquier fila con factibilidad $<$ 3.
\end{frame}

\begin{frame}{Reglas de decisión}
  \begin{itemize}
    \item Pedro puede vetar lo \textbf{genérico}.
    \item Mario puede vetar lo \textbf{mal definido} o salud.
    \item Ambición alta, pero ejecutable en 6 semanas.
    \item Solo \textbf{un} gancho de novedad (no inventar método + datos + framing a la vez).
    \item En el oral: Mario abre con datos/métricas; Pedro cierra con el porqué científico.
  \end{itemize}
\end{frame}

\begin{frame}{Siguiente paso concreto}
  \PipelineFive{Elegir\\principal}{Elegir\\reserva}{Bajar\\datos}{Explorar\\2--3 gráficos}{Escribir\\proposal}
  \vspace{0.5em}
  \begin{block}{La proposal (2 páginas) debe incluir}
    Pregunta precisa · motivación · dataset justificado · exploración preliminar · método inicial · plan realista.
  \end{block}
  \begin{center}\Large\textbf{¿Con cuál os quedáis?}\end{center}
\end{frame}

\end{document}
